\documentclass[11pt]{article}
\usepackage[margin=1.1in]{geometry}
\usepackage{amsmath,amssymb,amsthm}
\usepackage{booktabs}
\usepackage{microtype}
\usepackage[hidelinks]{hyperref}

\setlength{\parindent}{0pt}
\setlength{\parskip}{0.6\baselineskip}

\newtheorem{proposition}{Proposition}
\newtheorem{remark}{Remark}

\DeclareMathOperator{\LSE}{LSE}
\newcommand{\Prob}{\Pr}

\title{A likelihood model for local comparison of genome assemblies\\
{\large Model statement and derivations}}
\author{Anton Zamyatin}
\date{\today}

\begin{document}
\maketitle

\begin{abstract}
\noindent
We state a probabilistic model for adjudicating between two candidate genome
assemblies at a single locus using the long reads from which they were built.
The model has three components: a Poisson process governing where reads
originate, a normalised pair hidden Markov model giving the probability of an
observed read given its origin, and a recruitment procedure that defines the
locus as a closed comparison unit.  We show that the local comparison
agrees with the genome-wide one up to a computable tolerance, that the read
count must not be replaced by a locally profiled rate, and that the coverage rate must be
estimated from a source independent of the assemblies.  The reported statistic
is a robust score rather than a likelihood ratio, reducing to one under stated
ideal conditions, and its guarantees are statements about expectations whose
finite-depth error rates we quantify.  Two propositions establish that the more familiar conditional
formulation is provably blind to copy-number differences among identical
repeat units and to false duplication.
\end{abstract}

\tableofcontents

\section{Overview}

Two assemblies of the same individual, built from the same reads, disagree at a
locus.  We want to know which one the reads support, by how much, and with what
confidence.

The approach is to treat each candidate as a \emph{hypothesis about the source
of the reads} and ask which hypothesis assigns higher probability to what was
actually observed.  Three things must be specified to make this concrete:

\begin{enumerate}
    \item \textbf{Where reads come from.}  A sequencing run scatters read
          starts along the genome.  We model this as a Poisson process with a
          known rate.  Section~\ref{sec:origin}.
    \item \textbf{What a read looks like once it has an origin.}  The
          sequencer copies the template imperfectly.  We model this with a
          three-state pair HMM, completed into a proper distribution over
          observable reads.  Section~\ref{sec:emission}.
    \item \textbf{What ``the locus'' means.}  A local comparison is only
          well-posed if the reads and the sequence they touch form a closed
          set.  Section~\ref{sec:unit}.
\end{enumerate}

The reported score $\Delta S$ is built from a difference of log-likelihoods,
plus a deliberate robust departure for reads only one candidate can explain
(Section~\ref{sec:truncation}).  Everything else in this document is either a
derivation of that difference, an argument about how to estimate the one
external parameter it needs, or a statement of what the model cannot do.

Table~\ref{tab:notation} collects the notation.

\begin{table}[t]
\centering
\small
\begin{tabular}{ll}
\toprule
symbol & meaning \\
\midrule
$A_1, A_2$      & the two candidate assemblies \\
$L_1, L_2$      & total length of the \emph{recruited} sequence in each candidate (bases) \\
$k_1, k_2$      & number of disjoint recruited segments in each candidate \\
$h_{12}(p,Y)$   & retention indicator: 1 if a read with mark $Y$ from origin $p$ is recruited \\
$\alpha_A(p)$   & retention probability at origin $p$, $\int h_{12}(p,Y)e(Y\mid p,A)\,dY$ \\
$S^{(a)}_j, C^{(a)}_j$ & read $j$'s mass at origins inside / outside the unit, candidate $a$ \\
$L_A^{\mathrm{eff}}$ & effective exposure length under the recruitment rule \\
$R=\{r_1,\dots,r_n\}$ & the recruited read set \\
$n$             & number of recruited reads \\
$r_j$           & sequence of read $j$, of length $\ell_j$ \\
$Q_{ju}$        & Phred quality of base $u$ of read $j$ \\
$\bar\ell$      & mean read length \\
$p$             & an origin slot: a (position, strand) pair \\
$t = t_1\cdots t_m$ & template: assembly sequence at slot $p$, on the relevant strand \\
$Y=(r,Q)$       & a read mark: sequence and per-base qualities \\
$e_{jp}$        & $e(Y_j\mid p,A)$, the mark likelihood at origin $p$ \\
$Z(\ell,p,A)$   & normaliser over reads of length $\ell$ (Section~\ref{sec:normaliser}) \\
$\lambda$       & read starts per base per strand \\
$c$             & per-haplotype base coverage depth \\
$c_k$           & k-mer coverage read off the heterozygous peak \\
$\varepsilon_u$ & per-base error probability, $10^{-Q_u/10}$ \\
$s$             & log-probability cost of one mismatch, $\log\!\big(3(1-\varepsilon)/\varepsilon\big)$ \\
$\delta_I,\delta_D$ & gap-open probabilities (insertion, deletion) \\
$\gamma_I,\gamma_D$ & gap-extend probabilities \\
$\tau$          & termination probability \\
$\kappa_v$      & template sequence context at position $v$ \\
$m_1, m_2$      & counts of reads explicable by only one candidate \\
$C$             & cap on the evidence contributed by one such read (nats) \\
$\Delta\mathrm{LLR}$ & likelihood ratio over reads both candidates place \\
$\Delta S$      & the reported statistic; positive favours $A_1$ \\
$\eta$          & locality tolerance (Proposition~\ref{prop:locality}) \\
$q$             & probability a read from missing sequence is retained and uniquely placeable \\
\bottomrule
\end{tabular}
\caption{Notation used throughout.}
\label{tab:notation}
\end{table}

\section{The comparison unit}
\label{sec:unit}

A locus is defined by recursive read--segment recruitment.  Starting from a seed
interval, we alternately expand segments to the reads that align to them and
reads to the segments they touch, in \emph{both} assemblies, until the process
closes.  The result is a set of segments in each candidate, with total lengths
$L_1$ and $L_2$, and a single shared read set $R$.

Two properties of the closure matter for what follows.

\paragraph{Symmetry.}  Recruitment runs against both candidates and the read set
is their union.  Without this, reads recruited via $A_1$ would be selected for
their fit to $A_1$, and including them in the comparison would be circular.

\paragraph{Sealing.}  Expansion stops only at boundaries where the two
candidates carry identical sequence and no recruited read crosses outward with
appreciable probability.  This is what makes the local score meaningful, as
follows.

Write the whole-assembly length as $L_A^{\text{whole}} = L_A + L_{\text{out}}$
and split the read set into recruited and outside parts.  Because the two
candidates agree outside the unit, all terms involving $L_{\text{out}}$ cancel
in any difference.  Two residues do not cancel automatically, and both arise
the same way: a read may fit in more than one place.  A read we kept may also
fit outside the unit, so part of its explanation is invisible to the local
computation; and a read we did not keep may fit inside the unit, so it is
missing from the local computation altogether.  Both must be bounded rather
than assumed away.

\subsection{Recruited reads that also fit elsewhere}

For a recruited read $j$ and candidate $a$, split the origin sum into local and
outside parts,
\begin{equation}
    S^{(a)}_j = \!\!\sum_{p\in A^{(a)}_{\text{rec}}}\!\! e^{(a)}_{jp},
    \qquad
    C^{(a)}_j = \!\!\sum_{p\in A^{(a)}_{\text{out}}}\!\! e^{(a)}_{jp} .
\end{equation}
These sit inside a logarithm, so a nonzero $C^{(a)}_j$ does not cancel --- it
dilutes.  Crucially, $C^{(a)}_j$ cannot be zero: a pair HMM with positive
mismatch and indel probabilities assigns strictly positive mass to every
sufficiently long template window, so \emph{every} origin in the genome
contributes something.  The correct statement is therefore a bound, not an
identity.  The per-read discrepancy between the global and local comparison is
\begin{equation}
    \log\!\Big(1+\tfrac{C^{(1)}_j}{S^{(1)}_j}\Big)
    - \log\!\Big(1+\tfrac{C^{(2)}_j}{S^{(2)}_j}\Big).
\label{eq:per-read-discrepancy}
\end{equation}

\subsection{Unrecruited reads that also fit here}

Symmetrically, a read that was \emph{not} recruited still contributes to the
genome-wide comparison.  Writing $U_j$ for its mass at origins where the
candidates agree --- identical for both, by assumption --- and $T^{(a)}_j$ for
its mass at origins inside the unit of candidate $a$, its contribution to the
global difference is
\begin{equation}
    \log\!\Big(1+\tfrac{T^{(1)}_j}{U_j}\Big)
    - \log\!\Big(1+\tfrac{T^{(2)}_j}{U_j}\Big),
\label{eq:unrecruited-discrepancy}
\end{equation}
which the local computation omits entirely.  This term is nonzero whenever an
unrecruited read has candidate-specific placements inside the unit, which is
exactly the situation in segmental duplications and tandem arrays, and also for
reads excluded by the containment rule of Section~\ref{sec:exposure}.  It is
not, however, unbounded work: the reads concerned are those with a qualifying
alignment overlapping the unit that failed recruitment, a finite set
enumerable from the same alignment index.

\subsection{The bound}

\begin{proposition}[Locality to tolerance eta]
\label{prop:locality}
Suppose $A_1$ and $A_2$ carry identical sequence outside the recruited unit.
Then, for the \emph{uncapped} ratio $\Delta\mathrm{LLR}$ of
\eqref{eq:llr-common},
\begin{equation}
    \big|\Delta\mathrm{LLR}_{\text{genome}}
      - \Delta\mathrm{LLR}_{\text{local}}\big| \;\le\; \eta,
    \qquad
    \eta = \eta_{\text{in}} + \eta_{\text{out}}
           + \eta_{\text{exposure}} + \eta_{\text{tail}},
\end{equation}
\begin{align}
    \eta_{\text{in}} &= \sum_{j\,\text{recruited}}
      \left|\log\frac{S^{(1)}_j+C^{(1)}_j}{S^{(1)}_j}
          - \log\frac{S^{(2)}_j+C^{(2)}_j}{S^{(2)}_j}\right|, \\
    \eta_{\text{out}} &= \sum_{j\,\text{unrecruited}}
      \left|\log\frac{U_j+T^{(1)}_j}{U_j+T^{(2)}_j}\right|, \\
    \eta_{\text{exposure}} &= \lambda\,\Big|
        \textstyle\sum_p \alpha_{A_1}(p) - 2L_1^{\mathrm{eff}}
      - \sum_p \alpha_{A_2}(p) + 2L_2^{\mathrm{eff}}\Big| , \\
    \eta_{\text{tail}} &= \sum_{j,\,a}
        \log\!\Big(1+\tfrac{\text{omitted mass}^{(a)}_j}{S^{(a)}_j}\Big).
\end{align}
\end{proposition}

\begin{proof}
Apply the triangle inequality to the sum of
\eqref{eq:per-read-discrepancy} over recruited reads, of
\eqref{eq:unrecruited-discrepancy} over unrecruited reads, and of the exposure
and truncation discrepancies term by term; all remaining terms cancel by the
agreement of the candidates outside the unit.
\end{proof}

\begin{remark}[The bound must be taken on the difference]
\label{rem:tight-bound}
It is tempting to bound each candidate's term separately, giving
$\sum_j[\log(1+C^{(1)}_j/S^{(1)}_j) + \log(1+C^{(2)}_j/S^{(2)}_j)]$ and
similarly for $\eta_{\text{out}}$.  That bound is valid but vacuous.  The
dominant population of unrecruited reads with inside mass is not the
weakly-aligning one --- those carry little mass by construction --- but the
reads excluded by the containment rule of Section~\ref{sec:exposure}, which
have excellent alignments and are discarded on a purely geometric criterion.
At a $45$\,kb unit with $15$\,kb reads they are a third of all origins.  For
such a read $T_j$ is large and $U_j$ negligible, so the separated form reports
of order $10^4$ nats each and would declare every locus unsealed.  Yet its true
contribution is near zero, because a boundary-crossing read lies mostly in
shared flank sequence and is explained identically by both candidates:
$T^{(1)}_j \approx T^{(2)}_j$ and the difference cancels.  The reads that
genuinely contribute are those whose alignment reaches from the boundary into
\emph{disputed} sequence, and the form above counts exactly those.
\end{remark}

The last two components deserve naming.  The term $\eta_{\text{exposure}}$
measures the gap between the exact selected exposure $\lambda\sum_p\alpha_A(p)$
of \eqref{eq:alpha} and the closed form $2\lambda L_A^{\mathrm{eff}}$ used in
practice, which is exact only when the retention rule reduces to containment
alone.  The term $\eta_{\text{tail}}$ covers origins discarded below the
qualifying-placement threshold and probability mass leaving the dynamic
programming band; both are bounded by the retention threshold itself, so both
shrink as it is relaxed.

\begin{remark}[The tolerance does not cover $\Delta S$]
\label{rem:eta-scope}
Proposition~\ref{prop:locality} bounds the uncapped ratio.  It says nothing
about the capped term $C(m_1-m_2)$ of \eqref{eq:final}, which is not an
approximation error to be bounded but a deliberate departure from the
likelihood (Remark~\ref{rem:not-lr}).  A verdict is therefore local to
tolerance $\eta$ \emph{as a likelihood ratio}; as a reported score it inherits
in addition whatever the cap contributes, which is why $m_1$ and $m_2$ are
reported separately rather than folded into a single number.
\end{remark}

We say a locus is \emph{sealed to tolerance $\eta$}.  Both sums are computable
from an alignment index that stores every qualifying placement genome-wide, and
$\eta$ should be reported alongside every verdict: a difference of $10^3$ nats
at $\eta = 5$ is decisive, whereas the same difference at $\eta = 10^3$ is not
a local statement at all.  The corresponding additivity claim is likewise
approximate: over disjoint units,
\begin{equation}
    \Big|\Delta\mathrm{LLR}_{\text{genome}}
      - \sum_{\text{loci}} \Delta\mathrm{LLR}_{\text{local}}\Big|
    \;\le\; \sum_{\text{loci}} \eta .
\end{equation}

\subsection{Failure to close}

Recursive recruitment need not terminate at a local unit.  A repeat-connected
component can expand through every homologous copy in the genome, and imposing
a boundary at an arbitrary iteration produces a unit for which
Proposition~\ref{prop:locality} gives no useful tolerance.  The method must
therefore admit a third outcome besides a verdict for either candidate:
\emph{non-localisable}, declared when recruitment exceeds a stated iteration or
size limit.  Such loci are reported, not scored.

\section{Where reads come from}
\label{sec:origin}

Consider the assembly as a line of positions.  Sequencing scatters read-start
points along it.  Let
\begin{equation}
    \lambda = \text{expected number of read starts per base, per strand.}
\end{equation}
An assembly of length $L_A$ offers $2L_A$ slots, so the expected number of reads
it produces is $\mu = 2\lambda L_A$, and the observed number is Poisson:
\begin{equation}
    \Prob(N = n) = e^{-\mu}\frac{\mu^n}{n!}, \qquad \mu = 2\lambda L_A .
\end{equation}
Given that $n$ starts occurred, each is equally likely to sit in any slot:
\begin{equation}
    \Prob(\text{read } j \text{ originates at slot } p) = \frac{1}{2L_A}.
\end{equation}

Two features of this construction deserve emphasis, because the rest of the
document turns on them.  First, $L_A$ is a \emph{claim made by the assembly},
not an estimate of the truth: if the candidate asserts more sequence than
exists, it predicts more reads than were seen, and the likelihood falls
accordingly.  Second, $\lambda$ must be measured \emph{without reference to
either candidate} (Section~\ref{sec:lambda}); estimating it from the assembly
under test would make the comparison vacuous, as Proposition~\ref{prop:profile}
makes precise.

\section{The mark likelihood}
\label{sec:emission}

The quantity $e_{jp}$ used in Section~\ref{sec:likelihood} must be a
\emph{probability distribution over observable reads}, not merely an alignment
score.  This section constructs it.  Three things that a bare forward score
leaves undetermined are settled here: how much template an origin exposes, how
the process stops, and how read length is accounted for.

\subsection{The mark and its factorisation}

An observed read is a mark $Y = (r, Q)$ consisting of a base sequence
$r = r_1\cdots r_\ell$ and its per-base qualities $Q = Q_1\cdots Q_\ell$, with
$\ell = |r|$.  We factorise
\begin{equation}
    e(Y \mid p, A)
    \;=\; \underbrace{f(\ell)}_{\text{fragment length}}
    \;\cdot\; \underbrace{g(Q \mid \ell)}_{\text{quality profile}}
    \;\cdot\; \underbrace{\Prob(r \mid Q, \ell, p, A)}_{\text{sequence}} .
\label{eq:mark-factorisation}
\end{equation}
The first two factors describe the library and the instrument.  Neither depends
on the candidate, so both cancel from the observed-read terms of
$\Delta\mathrm{LLR}$ --- but not from the recruitment exposure, which is taken
as an integral against them (Remark~\ref{rem:library-factors}), so $f$ at least
must be estimated.  Read length is therefore \emph{modelled} rather than
ignored: it appears in $f$, and it is $f$ that supplies the mean $\bar\ell$ in
Proposition~\ref{prop:exposure}.  Qualities enter as conditioning covariates,
which is why the sequence term is written $\Prob(r\mid Q,\ell,p,A)$ and not
$\Prob(r\mid p,A)$; the model does not attempt to predict $Q$.

Only the third factor is candidate-dependent, and the remainder of this section
defines it.

\subsection{States and the template window}

Let the origin be $p = (s,o)$ and write $t = t_1 t_2 \cdots$ for the assembly
sequence read from position $s$ on strand $o$.  Walk along read and template
together.  At each step exactly one of three things happens, distinguished by
\emph{what is consumed}:

\begin{center}
\begin{tabular}{lll}
\toprule
state & consumes & interpretation \\
\midrule
$M$ & one read base and one template base & aligned column \\
$I$ & one read base only & insertion in the read \\
$D$ & one template base only & deletion from the read \\
\bottomrule
\end{tabular}
\end{center}

The number of template bases consumed, $m$, is \emph{not} a parameter supplied
with the origin: a start position does not determine how much template is
exposed once indels are permitted.  It is a latent variable, determined by the
path, and marginalised over.  This is the first thing that distinguishes a mark
likelihood from an alignment score.

There is no separate mismatch state.  A match and a mismatch consume exactly
the same thing and differ only in the probability of the letter emitted;
folding them into $M$ is exact unless substitution errors cluster, which they
do not appreciably do in HiFi data.  Indel errors \emph{do} cluster, and are
handled in Section~\ref{sec:context}.

\subsection{Transitions and termination}

\begin{center}
\begin{tabular}{lll}
\toprule
from $\to$ to & probability & meaning \\
\midrule
$B \to M$ & $1-\delta_I$ & begin aligned \\
$B \to I$ & $\delta_I$ & begin with an insertion \\
$M \to M$ & $1-\delta_I-\delta_D-\tau$ & keep matching \\
$M \to I$ & $\delta_I$ & open an insertion \\
$M \to D$ & $\delta_D$ & open a deletion \\
$M \to E$ & $\tau$ & terminate \\
$I \to I$ & $\gamma_I$ & extend insertion \\
$I \to M$ & $1-\gamma_I-\tau$ & close it \\
$I \to E$ & $\tau$ & terminate \\
$D \to D$ & $\gamma_D$ & extend deletion \\
$D \to M$ & $1-\gamma_D$ & close it \\
\bottomrule
\end{tabular}
\end{center}

The termination probability $\tau$ is what makes the model a distribution over
reads of \emph{variable} length; without an explicit $E$ state the forward sum
is a global-alignment score and nothing more.  Two transitions are absent by design, for the same reason.  Termination is
permitted from $M$ and $I$ but not from $D$: a deletion at the very end would
consume template beyond the last emitted read base, which no observation can
distinguish from terminating earlier.  Symmetrically, the begin state may enter
$M$ or $I$ but not $D$: a path beginning $B \to D$ consumes $t_1$ without
emitting, and is therefore identical to the path from origin $s+1$ beginning
$B \to M$.  Since the likelihood of Section~\ref{sec:likelihood} marginalises
over origins, admitting $B \to D$ would count the same alignment once at origin
$s$ and again at $s+1$.  Both exclusions remove unidentifiable redundancy
rather than real possibilities.

\begin{remark}
The standard pair HMM does admit $B \to D$, because it aligns two given
sequences from a fixed start and no redundancy arises.  The degeneracy here is
created by the marginalisation over origins, and is a consequence of embedding
the alignment model inside a point process rather than of the alignment model
itself.
\end{remark}

Separating gap opening ($\delta$) from gap extension ($\gamma$) makes indel
lengths geometric, which approximates the observed distribution far better than
a single per-base indel rate.  Note that $\gamma$ is used for extension here;
the symbol $\varepsilon$ is reserved throughout for per-base error
probabilities, and the two must not be conflated, since several claims in this
document turn on error calibration.

\subsection{Emissions}

With $\varepsilon_u = 10^{-Q_u/10}$ the error probability of read base $u$, the
aligned-state emission is
\begin{equation}
    \epsilon_M(u,v) =
    \begin{cases}
        1-\varepsilon_u & r_u = t_v,\\[2pt]
        \varepsilon_u/3 & r_u \neq t_v .
    \end{cases}
\end{equation}
State $D$ emits nothing and contributes a factor of one.  The cost of a single
mismatch follows:
\begin{equation}
    s = \log\frac{1-\varepsilon_u}{\varepsilon_u/3}
      = \log\frac{3(1-\varepsilon_u)}{\varepsilon_u} ,
\label{eq:mismatch-cost}
\end{equation}
which is $5.7$ nats at $Q20$, $8.0$ at $Q30$ and $10.3$ at $Q40$.  Using the
read's own qualities rather than a flat rate is not a refinement: in a repeat
array where a handful of variant columns separate the candidates, it is the
difference between a decisive and an ambiguous read.

The insertion emission requires a choice.  Taking $\epsilon_I(u) = 1/4$ treats
an inserted base as uniformly random and ignores both its reported quality and
its context.  For HiFi this is demonstrably crude: inserted bases occur
overwhelmingly within homopolymer runs and short tandem repeats, where the
inserted base is very nearly determined by the surrounding template.  We
therefore allow a context-conditioned insertion emission
$\epsilon_I(u \mid \kappa_v)$, estimated empirically alongside the gap
parameters, and reduce to $1/4$ only outside repetitive context.  Whichever is
used must be stated and calibrated on held-out data, since it enters every
read that crosses a homopolymer.

\subsection{The forward recursion}

A particular walk through the states is an alignment path $\pi$, whose
probability is the product of the transitions and emissions along it.  Many
paths produce the same read from the same template, and which one occurred is
unobserved, so the paths must be summed.  Define
\begin{equation}
    f_X(u,v) = \Prob\big(r_1..r_u \text{ emitted},\ t_1..t_v \text{ consumed},\
    \text{currently in state } X\big),
    \qquad X \in \{M,I,D\},
\end{equation}
each already a sum over all paths reaching that configuration.  Then
\begin{align}
    f_M(u,v) &= \epsilon_M(u,v)\Big[
        f_M(u{-}1,v{-}1)(1-\delta_I-\delta_D-\tau) \notag \\
      &\qquad\qquad
      + f_I(u{-}1,v{-}1)(1-\gamma_I-\tau)
      + f_D(u{-}1,v{-}1)(1-\gamma_D)\Big], \\
    f_I(u,v) &= \epsilon_I(u)\Big[f_M(u{-}1,v)\,\delta_I + f_I(u{-}1,v)\,\gamma_I\Big], \\
    f_D(u,v) &= f_M(u,v{-}1)\,\delta_D + f_D(u,v{-}1)\,\gamma_D,
\end{align}
initialised from the begin state at $(0,0)$ with the $B$ transitions above, so
that $f_D(0,v) = 0$ for all $v$.
Terminating and marginalising over the template window gives
\begin{equation}
    \Prob(r,\ell \mid Q, p, A)
    = \tau \sum_{m} \big[f_M(\ell,m) + f_I(\ell,m)\big].
\label{eq:joint-r-ell}
\end{equation}
The sum over $m$ is the marginalisation the fixed-window formulation omits.

\subsection{The normaliser}
\label{sec:normaliser}

Equation~\eqref{eq:joint-r-ell} is a joint distribution over sequence and
length, whose implied length distribution is approximately geometric with mean
$1/\tau$ --- not the length distribution of a size-selected HiFi library.  We
therefore condition on the observed length and let the empirical $f(\ell)$ of
\eqref{eq:mark-factorisation} carry it:
\begin{equation}
    \Prob(r \mid Q,\ell,p,A) = \frac{\Prob(r,\ell\mid Q,p,A)}{Z(\ell,p,A)},
    \qquad
    Z(\ell,p,A) = \sum_{r'} \Prob(r',\ell\mid Q,p,A).
\end{equation}
The normaliser is not an extra approximation and needs no separate machinery:
for any fixed path the emissions sum to one over $r'$, so $Z$ is the total
transition mass of paths emitting exactly $\ell$ bases, obtained by running the
same recursion with every emission probability set to one.

If the gap parameters are context-independent, $Z(\ell,p,A)$ does not depend on
$p$ or $A$ and cancels in $\Delta\mathrm{LLR}$; it may then be omitted
entirely.  With the context-dependent parameters of Section~\ref{sec:context}
it does not cancel, and the origin sum
\begin{equation}
    \sum_p e_{jp}
    = \sum_p \frac{\Prob(r_j,\ell_j\mid Q_j,p,A)}{Z(\ell_j,p,A)}
\label{eq:origin-sum-normalised}
\end{equation}
is a sum of \emph{ratios}, with a different denominator at each origin.  This
has an algorithmic consequence, taken up in Section~\ref{sec:computation}.

\begin{remark}[Where the template ends]
The construction above lets the walk consume template indefinitely.  Within a
recruited segment the template is finite, and paths reaching a segment boundary
carry mass that leaves the unit.  This is not a separate defect: it is exactly
the containment thinning of Section~\ref{sec:exposure}, whose exposure
consequences are given by Proposition~\ref{prop:exposure} and whose residual
placement mass is accounted in the tolerance $\eta$ of
Proposition~\ref{prop:locality}.
\end{remark}

\begin{remark}[Why the sum, not the best path]
\label{rem:viterbi}
Take template \texttt{AAAAAAA} and read \texttt{AAAAAA}.  One base was deleted,
but seven distinct paths place that deletion, all with identical probability
$q$.  Viterbi reports $q$; the forward sum reports $7q$, a difference of
$\log 7 = 1.95$ nats --- and the forward value is the correct one, because the
read really is seven times more likely than any single path suggests.  In
homopolymer-rich sequence this discrepancy is large, systematic, and
\emph{different between the two candidates}, since they have different run
lengths.  Viterbi therefore injects a bias directly into the quantity being
measured.  This is the reason the aligner's CIGAR must not be reused: alignment
is used to \emph{locate} candidate origins, never to score them.
\end{remark}

\subsection{Context-dependent gap parameters}
\label{sec:context}

A single set $(\delta_I,\delta_D,\gamma_I,\gamma_D)$ is misspecified exactly
where HiFi errors concentrate.  We therefore let the gap parameters depend on
the local template context $\kappa_v$ --- minimally the homopolymer run length
and its base, optionally the period of a short tandem repeat:
\begin{equation}
    \delta_I(\kappa_v),\quad \delta_D(\kappa_v),\quad
    \gamma_I(\kappa_v),\quad \gamma_D(\kappa_v).
\end{equation}
For run lengths $1$--$10$ over four bases this is roughly forty parameters,
estimated once from reads aligned to a validated reference, from undisputed
sequence, and then held fixed across both candidates.  Fitting them per
candidate would let a poorer assembly compensate with looser error rates.  The
motivation is empirical: homopolymer richness is the strongest predictor of
anomalous HiFi behaviour in the loci this method targets, so a
context-independent model is calibrated everywhere except where it is needed.

\subsection{Computation}
\label{sec:computation}

Probabilities underflow double precision within a few hundred bases, so all
recursions are carried in logarithms with
$\LSE(a,b) = \max(a,b) + \log(1+e^{-|a-b|})$.

Origins must not be enumerated one at a time.  Running a banded forward pass at
every candidate origin costs $O(L\ell w)$ per read for a unit of length $L$,
roughly $6\times10^{11}$ cells for a $15$\,kb read against a $100$\,kb unit,
which is not affordable.  A single dynamic program over the whole segment with
a free start state sums over origins internally at $O(L\ell)$, two orders of
magnitude cheaper.

That free-start pass computes $\sum_p \Prob(r_j,\ell_j\mid Q_j,p,A)$: the sum
of the \emph{numerators} of \eqref{eq:origin-sum-normalised}.  When $Z$ is
origin-dependent this is not the quantity required, and no single pass over the
segment produces a sum of origin-specific ratios.  We therefore use a two-tier
scheme:

\begin{enumerate}
    \item a free-start pass over the segment, at $O(L\ell)$, to \emph{locate}
          the origins carrying non-negligible numerator mass;
    \item per-origin normalisation on that retained set only, at
          $O(k\ell w)$ for $k$ retained origins, evaluating both
          $\Prob(r_j,\ell_j\mid Q_j,p,A)$ and $Z(\ell_j,p,A)$ by the recursion
          of Section~\ref{sec:normaliser}.
\end{enumerate}

Because $Z$ varies far more slowly across origins than the numerator does,
tier one is a safe locator even though it uses unnormalised mass; the residual
is bounded by the retention threshold and enters $\eta_{\text{tail}}$.  With
context-independent gaps the scheme collapses to tier one alone.  In neither
case is the origin set capped by an aligner's secondary-alignment limits, which
matters because that sum carries the copy-number signal.

Banding remains necessary in the read direction.  A fixed diagonal band is not
automatically safe in indel-rich repeats, so the band should be adaptive and
the probability mass reaching its edge should be reported as a diagnostic,
in the same spirit as the tolerance $\eta$.

\section{The likelihood}
\label{sec:likelihood}

Combining the count, the origin prior and the emission model, the probability of
observing read set $R$ of size $n$ is
\begin{equation}
    \Prob(R,n \mid A)
    = \underbrace{e^{-2\lambda L_A}\frac{(2\lambda L_A)^n}{n!}}_{\text{how many reads}}
      \times
      \underbrace{\prod_{j=1}^{n}\Big[\frac{1}{2L_A}\sum_{p} e_{jp}\Big]}_{\text{what they look like}} .
\end{equation}
Taking logarithms, the factors $(2L_A)^n$ and $(2L_A)^{-n}$ cancel and all
length dependence collapses into a single term:
\begin{equation}
    \boxed{\;
    \log\Prob(R,n\mid A)
    = -2\lambda L_A \;+\; n\log\lambda \;-\; \log n!
      \;+\; \sum_{j=1}^{n}\log\sum_{p}e_{jp} \; }
\label{eq:loglik}
\end{equation}

The three groups of terms read as follows.  The term $-2\lambda L_A$ is a rent
paid on every base the assembly claims: more sequence means more reads should
have been seen.  The middle terms contain no reference to $A$ and are constants
of the comparison.  The final sum measures how well the reads are explained.

Comparing two candidates, the constants vanish and the origin slot sets are of
course different sizes:
\begin{equation}
    \Delta\mathrm{LLR}
    = \sum_{j=1}^{n}\Bigg[
        \log\!\!\sum_{p\in\mathrm{slots}(A_1)}\!\! e^{(1)}_{jp}
      - \log\!\!\sum_{p\in\mathrm{slots}(A_2)}\!\! e^{(2)}_{jp}
      \Bigg]
      \;-\; 2\lambda\,(L_1-L_2).
\label{eq:dllr}
\end{equation}
Positive values favour $A_1$.  The inner sums are \emph{totals}, not
distributions --- the normalisation was already extracted into the length term
--- so a read with several plausible origins genuinely scores higher than one
with a single origin.  That multiplicity is the model's copy-number signal.

\begin{remark}[Units of the penalty]
Writing $\lambda = c/(2\bar\ell)$ gives
$2\lambda\,\Delta L = c\,\Delta L/\bar\ell$, which is exactly the number of
reads expected to start within $\Delta L$ extra bases.  The length penalty is
denominated in \emph{missing reads}.  At $30\times$ with $15$\,kb reads, ten
kilobases of unsupported sequence costs twenty nats, being the twenty reads that
should have been there.
\end{remark}

\section{Exposure under recruitment}
\label{sec:exposure}

Equation~\eqref{eq:loglik} treats $R$ as the complete set of reads originating
in the candidate, with exposure $2\lambda L_A$.  The recruited set is not that:
it is produced by a rule applied to alignments, and that rule discards reads.
Since the rule is applied to each candidate, the discarding is in principle
candidate-dependent and cannot be assumed ignorable.  This section states the
exposure exactly, and shows which part of it cancels.

\subsection{The selected process}

Let $Y = (r,Q)$ denote a read mark, with $\ell = |r|$.  Recruitment is a
\emph{deterministic} rule applied to the frozen pair of candidates: write
\begin{equation}
    h_{12}(p,Y) \in \{0,1\}
\end{equation}
for the indicator that a read with mark $Y$ arising at origin $p$ is retained.
Two features of this definition are deliberate.  It is an indicator, not a
probability: given the frozen segments and a fixed threshold, retention is a
function of the data, and randomness enters only through $Y$.  And it is
indexed by the \emph{pair} $(A_1,A_2)$, not by a single candidate: under union
recruitment a read may be retained through its alignment to the other
candidate, which a per-candidate rule cannot express.  The same $h_{12}$ is
therefore applied when scoring against either candidate, which is what makes
the comparison symmetric.

Retained reads form a thinned marked Poisson process with intensity
$\lambda\,h_{12}(p,Y)\,e(Y\mid p,A)$, whose log-likelihood under candidate $A$
is
\begin{equation}
    \ell_A
    = -\lambda\!\!\sum_{p\in\mathrm{slots}(A)}\!\! \alpha_A(p)
      \;-\;\log n!
      \;+\; \sum_{j=1}^{n}\log\Big[\lambda\!\!\sum_{p\in\mathrm{slots}(A)}\!\!
              h_{12}(p,Y_j)\,e(Y_j\mid p,A)\Big],
\label{eq:selected-loglik}
\end{equation}
\begin{equation}
    \alpha_A(p) = \int h_{12}(p,Y)\,e(Y\mid p,A)\,dY .
\label{eq:alpha}
\end{equation}
Two changes relative to \eqref{eq:loglik} are visible.  The exposure is
$\lambda\sum_p \alpha_A(p)$ rather than $2\lambda L_A$; and the origin sum is
weighted by $h_{12}$, so it runs only over origins at which the observed read
would have been retained.  The second is automatic in any implementation that
scores only admitted origins.  The first requires evaluation, and is the
subject of Proposition~\ref{prop:exposure}.

\begin{remark}[Which library factors are needed]
\label{rem:library-factors}
The factors $f(\ell)$ and $g(Q\mid\ell)$ of \eqref{eq:mark-factorisation}
cancel from the observed-read terms of \eqref{eq:selected-loglik}, since they
do not depend on the candidate.  They do \emph{not} cancel from
\eqref{eq:alpha}: retention depends on read length through containment, and on
quality through the placeability threshold, so the integral defining
$\alpha_A(p)$ is taken against them.  An empirical fragment-length
distribution is therefore required, and the quality profile is required
whenever the phase-two rule is quality-dependent.  This is the sense in which
the length distribution is part of the model rather than a nuisance.
\end{remark}

\begin{remark}[Phase one is still post-selection]
\label{rem:phase-one-selection}
Freezing the segments before scoring removes the dependence of $h_{12}$ on the
expansion path, but not the dependence of the \emph{unit itself} on the reads
that later score it.  The rule is symmetric in the candidates by construction,
so it does not obviously favour either; what it can bias is the unit
definition, hence $L_1 - L_2$.  Ignorability is not established, and the clean
remedy is the cross-fitting of Section~\ref{sec:validation}: recruit and freeze
on one fold of molecules, score on another.
\end{remark}

\subsection{Two-phase recruitment}
\label{sec:two-phase}

We separate the rule into two phases, because they play different roles and
carry different risks.

\emph{Phase one} determines the comparison unit.  Reads and segments are
expanded recursively under a strict identity and coverage-fraction criterion,
and iteration continues to convergence.  The criterion must be strict: a
permissive rule admits paralogous sequence, which recruits its own reads, and
the unit fails to close.  The output is a frozen set of $k_1$ and $k_2$ disjoint
segments with total lengths $L_1$ and $L_2$.

\emph{Phase two} determines which reads are scored.  Over the frozen segments,
and without any further segment expansion, reads are collected under a second,
permissive rule and each is scored against both candidates.  Because no
expansion occurs, a permissive rule here cannot enlarge the unit; a read that
in truth belongs elsewhere simply scores poorly against both candidates and
contributes negligibly to the difference.

The separation matters for the model as well as for the implementation: after
phase one the segments are fixed, so $a_A(p,Y)$ depends only on the frozen
segments and the phase-two rule, and not on the order in which recruitment
proceeded.

\subsection{Exposure under containment}

The phase-two rule retains a read only if its alignment falls entirely within
one segment.  Reads originating near a segment end therefore cannot be
retained, and the loss is substantial at the scale of a locus: for a segment of
$45$\,kb and a mean read length of $15$\,kb, a third of all origins are
excluded.  Nevertheless:

\begin{proposition}[Containment exposure]
\label{prop:exposure}
Let candidate $A$ consist of $k_A$ disjoint segments of lengths
$L_A^{(1)},\dots,L_A^{(k_A)}$ summing to $L_A$, and let read reference spans be
drawn from a distribution $f$ with mean $\bar\ell$, independent of the
candidate.  Under the containment rule,
\begin{equation}
    \sum_{p}\alpha_A(p)
    = 2\sum_{i=1}^{k_A}\sum_{\ell} f(\ell)\,\max\!\big(0,\;L_A^{(i)}-\ell+1\big)
    \;\equiv\; 2 L_A^{\mathrm{eff}} .
\end{equation}
If every segment is longer than the longest read, this reduces to
\begin{equation}
    L_A^{\mathrm{eff}} = L_A - k_A(\bar\ell - 1),
\label{eq:leff}
\end{equation}
and the exposure difference between candidates is
\begin{equation}
    -2\lambda\big(L_1^{\mathrm{eff}} - L_2^{\mathrm{eff}}\big)
    = -2\lambda\Big[(L_1-L_2) - (k_1-k_2)(\bar\ell-1)\Big].
\label{eq:exposure-diff}
\end{equation}
\end{proposition}

\begin{proof}
A read of reference span $\ell$ is contained in segment $i$ precisely when its
start lies in $[\,\text{start}_i,\ \text{end}_i-\ell\,]$, giving
$\max(0, L_A^{(i)}-\ell+1)$ admissible starts per strand.  Averaging over $f$
and summing over segments and strands gives the first display.  When
$L_A^{(i)} > \ell$ for all $i$ and $\ell$ in the support of $f$, the maximum is
inactive and the inner sum is $L_A^{(i)} - \bar\ell + 1$; summing over $i$ gives
\eqref{eq:leff}.  Subtracting the two candidates gives
\eqref{eq:exposure-diff}.
\end{proof}

Two consequences.

\paragraph{The boundary loss is common mode.}  The term $\bar\ell$ appears with
the same coefficient in both candidates and cancels, because both are scored
against the same reads and hence the same length distribution.  The
containment rule therefore discards a third of the data at a $45$\,kb locus
without perturbing the length penalty at all, provided the segment counts agree.
This is what licenses the use of $-2\lambda(L_1-L_2)$ in
\eqref{eq:dllr}, which would otherwise be an unexamined assumption.

\paragraph{Fragmentation does not cancel.}  When $k_1 \neq k_2$ the residual
term $2\lambda(k_1-k_2)(\bar\ell-1)$ survives, favouring the candidate split
into more segments.  At $30\times$ coverage with $15$\,kb reads this is
approximately $30$ nats per additional segment, which is comparable to the
signal from a genuine event and points in a systematic direction, since
differing fragmentation is a common form of disagreement at difficult loci.
The correction in \eqref{eq:exposure-diff} removes it and costs nothing to
compute.

\begin{remark}
Equation~\eqref{eq:leff} requires every segment to exceed the longest read.
Implementations should verify this and otherwise evaluate the exact form,
$\sum_{\ell\le L}f(\ell)(L-\ell+1) = (L+1)F(L) - \sum_{\ell\le L}\ell f(\ell)$,
which remains a difference of candidate-specific constants and preserves the
argument.  Containment is a condition on reference coordinates, so $\ell$ must be the
\emph{reference span}, not the query length.  Using the observed read-length
distribution in its place is an approximation: the two differ by the net indel
count, which is context-dependent and, in indel-rich sequence, mildly
candidate-dependent.  The residual belongs to $\eta_{\text{exposure}}$ and
should be estimated from the span distribution directly rather than assumed
negligible.
\end{remark}

\begin{remark}[What is not claimed]
Proposition~\ref{prop:exposure} addresses only the containment component of
$a_A(p,Y)$.  The phase-two identity criterion also thins, and does so in a
candidate-dependent way: reads spanning a structural difference fail against one
candidate and not the other.  That thinning is not a nuisance to be cancelled
--- it is the definition of the classes $m_1$ and $m_2$ of
Section~\ref{sec:truncation}, and it must be pre-registered and swept rather
than tuned.
\end{remark}

\section{Why local coverage must not be profiled away}
\label{sec:conditioning}

A widespread alternative treats $n$ as fixed and scores only the placement of
reads, giving a penalty $-n\log(2L_A)$.  The objection is narrower than a
prohibition on conditioning, and worth stating precisely:

\begin{quote}
Do not estimate a separate coverage rate \emph{within each candidate locus},
because doing so profiles away the exposure contrast that identifies copy
number.
\end{quote}

Conditioning as such is not the problem.  Conditioning on a shared global
yield, and modelling the allocation of read starts between the locus and
external control sequence as multinomial with probabilities proportional to
candidate exposures, preserves the contrast perfectly well.  What is required
is that the locus be compared against an exposure it does not itself set.  The
formulation below achieves that by fixing $\lambda$ externally; the multinomial
route achieves it by comparison to a control region.  The failure mode is
specific and is what the following establishes.

\begin{proposition}[The conditional model is the profiled Poisson model]
\label{prop:profile}
Maximising \eqref{eq:loglik} over $\lambda$ within the unit gives
$\hat\lambda = n/(2L_A)$; substituting back yields
\begin{equation}
    -2\hat\lambda L_A + n\log\hat\lambda
    = -n + n\log n - n\log(2L_A),
\end{equation}
whose only assembly-dependent term is $-n\log(2L_A)$.  The conditional
likelihood is therefore the profile likelihood of the Poisson model with the
rate estimated locally.
\end{proposition}

The consequence is that if $\lambda$ is free to adapt to the disputed region, a
collapsed assembly simply reports higher coverage and nothing is learned.  All
information about copy number lies in the statement \emph{more reads landed here
than this much sequence should attract}, and that statement is empty unless
$\lambda$ is fixed externally.  Two concrete failures follow.

The two propositions that follow are stated in an \emph{ideal regime}: no
selection, no containment, and no locality error, so that exposure is $2\lambda
L_A$ with the raw length rather than $2\lambda L_A^{\mathrm{eff}}$.  This
isolates the effect of conditioning from the corrections of
Section~\ref{sec:exposure}; substituting $L^{\mathrm{eff}}$ throughout changes
the constants but not the conclusions, since the correction is common to both
formulations.

\begin{proposition}[Blindness to tandem copy number]
\label{prop:array}
Let a comparison unit consist of unique flank of length $U$ plus an array of
identical copies of a unit of length $\rho$.  Let $A_1$ carry $k_1$ copies and
$A_2$ carry $k_2 < k_1$, both with well-formed junctions, so that every read is
placeable in both candidates.  Reads inside the array then have $k_1$ plausible
origins in $A_1$ and $k_2$ in $A_2$, and
\begin{align}
    \Delta\mathrm{LLR}_{\text{conditional}}
      &= n\,\frac{k_1\rho}{L_1}\log\frac{k_1}{k_2}
       - n\log\frac{L_1}{L_2},\\
    \mathbb{E}\big[\Delta\mathrm{LLR}_{\text{Poisson}}\big]
      &= 2\lambda\rho\Big(k_1\log\frac{k_1}{k_2} - (k_1-k_2)\Big).
\end{align}
When the unit is the array alone ($U=0$) the conditional expression is
identically zero for every $k_1, k_2, \rho$ and every coverage, whereas the
Poisson expression is strictly positive and independent of $U$.
\end{proposition}

\begin{proof}
At $U=0$ we have $L_1 = k_1\rho$, $L_2 = k_2\rho$ and $k_1\rho/L_1 = 1$, so the
conditional expression is $n\log(k_1/k_2) - n\log(k_1/k_2) = 0$.  For the
Poisson form, the expected number of reads originating in the array is
$2\lambda k_1\rho$, each gaining $\log(k_1/k_2)$, against a length penalty
$2\lambda(k_1-k_2)\rho$.  Positivity follows from
$k_1\log(k_1/k_2) > k_1-k_2$, which is the strict concavity of the logarithm
applied to $\log(k_2/k_1) < k_2/k_1 - 1$.
\end{proof}

The conditional model is thus not merely biased toward collapses: it is exactly
blind to them, because the multiplicity gain $\log(k_1/k_2)$ and the parsimony
penalty $-\log(L_1/L_2)$ are the same quantity with opposite signs.  Note the
perverse interaction with recruitment: tighter closure means smaller $U$, hence
less surviving signal, whereas under the Poisson formulation the verdict does
not depend on how much flank the unit happens to contain.

\begin{table}[t]
\centering
\small
\begin{tabular}{lrrrr}
\toprule
$k_1 \to k_2$ & explanation & penalty & Poisson & conditional \\
\midrule
$5 \to 3$  & 229.9 & 180.0 & \textbf{49.9}  & 0 \\
$5 \to 4$  & 100.4 &  90.0 & \textbf{10.4}  & 0 \\
$10 \to 5$ & 623.8 & 450.0 & \textbf{173.8} & 0 \\
$2 \to 1$  & 124.8 &  90.0 & \textbf{34.8}  & 0 \\
\bottomrule
\end{tabular}
\caption{Expected $\Delta\mathrm{LLR}$ in nats for partial collapse of a
tandem array, at $\rho = 45$\,kb, $\bar\ell = 15$\,kb, $30\times$, $U=0$.}
\label{tab:collapse}
\end{table}

\begin{remark}[Junction reads, and why the proposition assumes them away]
\label{rem:junction}
Proposition~\ref{prop:array} assumes both candidates carry well-formed
junctions.  That assumption is doing real work, and dropping it changes the
picture.

Consider instead a \emph{total} collapse, $k_2 = 1$, in which $A_2$ retains one
copy followed directly by flank.  The $k_1-1$ internal junctions
$\ldots\rho_{\mathrm{end}}\mid\rho_{\mathrm{start}}\ldots$ of $A_1$ then exist
nowhere in $A_2$: a read straddling one is a circular permutation of $\rho$ and
matches no linear substring of the collapsed candidate.  Such reads are not
multiplicity reads at all but members of the one-sided class $m_1$, and their
expected number is $2\lambda(k_1-1)(\bar\ell-1)$.  At $k_1=2$,
$\rho=45$\,kb and $30\times$ that is $30$ reads against $150$ interior ones,
contributing $C m_1 = 3000$ nats at $C=100$ --- two orders of magnitude above
the $34.8$ nats of the multiplicity channel.  The conditional model sees this
signal too, since it does not pass through the multiplicity term at all.

The proposition is therefore not a claim that collapsed arrays are undetectable
under conditioning.  It is the narrower and more useful claim that the
\emph{copy-number channel} is dead under conditioning, and that channel is all
that remains when an assembler produces a structurally plausible partial
collapse --- which is the common case, and the one this method exists to
adjudicate.

The same example gives a closed form for the detection probability $q$ of
Proposition~\ref{prop:count}.  Reads originating in the missing $(k_1-1)\rho$
of sequence are not all detectable: interior ones are explained perfectly by
the surviving copy.  Only junction readers are one-sided, so
\begin{equation}
    q = \frac{\bar\ell-1}{\rho},
\end{equation}
which is $1/3$ in this example.  Accordingly $m_1/(2\lambda) = 15$\,kb
recovers $(k_1-1)\bar\ell$ and understates the missing length, whereas
$m_1/(2\lambda q) = 45$\,kb recovers $(k_1-1)\rho$ exactly.  When reads exceed
the unit length every array read spans a junction and $q \to 1$.
\end{remark}

The Poisson expression in Proposition~\ref{prop:array} is an expectation.  The
realised score is random, because the number of reads originating in the array
is $N \sim \text{Poisson}(2\lambda k_1\rho)$ while the length penalty is fixed:
\begin{equation}
    \Delta = N\log\frac{k_1}{k_2} - 2\lambda(k_1-k_2)\rho,
    \quad
    \operatorname{Var}[\Delta] = 2\lambda k_1\rho\Big(\log\frac{k_1}{k_2}\Big)^2 .
\label{eq:array-moments}
\end{equation}
The verdict is wrong whenever $N \le 2\lambda(k_1-k_2)\rho/\log(k_1/k_2)$,
which at low depth is not a remote event.  Table~\ref{tab:power} gives the
failure probabilities for the $2\to1$ case.

\begin{proposition}[Blindness to false duplication]
\label{prop:dup}
Let a region of true length $\ell_{\mathrm{loc}}$ be correctly represented by
$A_1$ ($L_1 = \ell_{\mathrm{loc}}$) and falsely duplicated in $A_2$
($L_2 = 2\ell_{\mathrm{loc}}$), with all $n = 2\lambda\ell_{\mathrm{loc}}$
observed reads originating from the region.  Every read acquires a second
identical origin in $A_2$, contributing $-n\log 2$ to the emission difference,
while the length term contributes
$-2\lambda(L_1-L_2) = 2\lambda\ell_{\mathrm{loc}} = n$.  Then
\begin{equation}
    \Delta\mathrm{LLR}_{\text{Poisson}} = n\big(1-\log 2\big) \approx 0.307\,n > 0,
\end{equation}
whereas the conditional version replaces the length term by
$-n\log(L_1/L_2) = n\log 2$ and gives $n\log 2 - n\log 2 = 0$.
\end{proposition}

Again this is an expectation.  With $N \sim \text{Poisson}(2\lambda
\ell_{\mathrm{loc}})$ the realised score is
$\Delta = 2\lambda\ell_{\mathrm{loc}} - N\log 2$, with mean
$2\lambda\ell_{\mathrm{loc}}(1-\log 2)$ and variance
$2\lambda\ell_{\mathrm{loc}}(\log 2)^2$; an unusually well-covered false
duplicate can invert the sign.

\begin{table}[t]
\centering
\small
\begin{tabular}{lrrrr}
\toprule
coverage & $1\times$ & $3\times$ & $10\times$ & $30\times$ \\
\midrule
collapsed array ($k=2$, $\rho=45$\,kb) & 0.29 & 0.092 & 0.013 & $3.9\times10^{-5}$ \\
false duplication ($\ell_{\mathrm{loc}}=45$\,kb) & 0.19 & 0.124 & 0.010 & $4.4\times10^{-5}$ \\
haplotype loss, $d=0$ ($\ell_{\mathrm{loc}}=45$\,kb) & 0.29 & 0.092 & 0.013 & $3.9\times10^{-5}$ \\
\bottomrule
\end{tabular}
\caption{Probability that the realised verdict has the wrong sign, at
$\bar\ell=15$\,kb, under the ideal assumptions of the corresponding
proposition.  These are floors: they exclude misspecification, selection, and
locality error.  The propositions establish the sign of
$\mathbb{E}[\Delta]$, not of every realisation.}
\label{tab:power}
\end{table}

Table~\ref{tab:power} is a power calculation rather than a caveat: it states
the depth required for a verdict at a given event size, and it is the right
form in which to report the method's resolution.

\section{The coverage rate}
\label{sec:lambda}

By Proposition~\ref{prop:profile}, $\lambda$ must come from outside the
assemblies.  The k-mer spectrum of the reads provides it, using no assembly at
any stage.

\paragraph{Step 1: read the peak.}  In a diploid spectrum, heterozygous k-mers
(present on one haplotype) peak at multiplicity $c_k$ and homozygous k-mers at
$2c_k$.  Because the assembly represents both haplotypes as separate sequence,
$L_A$ counts both, so $\lambda$ must be the \emph{per-haplotype} rate: the
first peak.  Where heterozygosity is too low for a resolved het peak, halve the
homozygous peak.  Fitting should be done with a spectrum model
(e.g.\ GenomeScope2) rather than by eye.

\paragraph{Step 2: correct for k-mers destroyed by errors.}  A k-mer is counted
only if all $k$ bases are correct, so
\begin{equation}
    c_k^{\text{true}} = \frac{c_k}{(1-\varepsilon)^k},
\end{equation}
a $2\%$ correction at $k=21$, $\varepsilon = 10^{-3}$.

\paragraph{Step 3: k-mer coverage to base coverage.}  A read of length $\ell$
carries $\ell-k+1$ k-mers but $\ell$ bases:
\begin{equation}
    c = c_k^{\text{true}}\cdot\frac{\bar\ell}{\bar\ell-k+1},
\end{equation}
a factor of $1.0013$ for HiFi at $k=21$.

\paragraph{Step 4: base coverage to read-start rate.}  Total sequenced bases are
$n\bar\ell$ over a haplotype of length $G$, so $c = n\bar\ell/G$, while
$n = 2\lambda G$ by definition.  Hence
\begin{equation}
    \boxed{\;\lambda = \frac{c}{2\bar\ell}\;}
\end{equation}
At $c=30\times$ and $\bar\ell=15$\,kb this gives $\lambda = 10^{-3}$: one read
start per kilobase per strand.

\subsection{Required accuracy}
\label{sec:required-accuracy}

$\lambda$ enters only through $-2\lambda(L_1-L_2)$.  Where sequence evidence is
strong the verdict is far from a tie and a few percent of error in $\lambda$
changes confidence, not direction.  The exception is the low-divergence regime.
Taking the haplotype comparison of Section~\ref{sec:diploid} at zero divergence,
$\Delta\mathrm{LLR} = n\log 2 - 2\hat\lambda\ell_{\mathrm{loc}}$ with
$n = 4\lambda\ell_{\mathrm{loc}}$, so the verdict inverts when
\begin{equation}
    \frac{\hat\lambda}{\lambda} > 2\log 2 \approx 1.386 .
\end{equation}
Roughly forty percent of headroom: ordinary estimation error is safe, but
choosing the wrong spectrum peak is a factor of two and lands squarely on the
wrong side.  That single choice matters more than every other approximation in
this document.

\subsection{Assumptions behind the spectrum estimate}

The translation above is first-order and rests on conditions that should be
stated rather than assumed: that k-mers are counted canonically; that the
ploidy peaks are resolved well enough for the heterozygous component to be
identified rather than inferred; that copy-number structure in the genome does
not distort the fitted peak, which it does when a large fraction of the
sequence is high-copy repeat; and that the error correction of Step~2 uses a
per-base error rate representative of the library rather than a nominal one.
Each is checkable, and each fails in a direction that biases $\lambda$.

\subsection{\texorpdfstring{Propagating uncertainty in $\lambda$}{Propagating uncertainty in lambda}}

The spectrum supplies an \emph{estimate}, not a constant, and
Section~\ref{sec:required-accuracy} shows the verdict is sensitive to it in the
low-divergence regime.  Treating $\hat\lambda$ as known therefore understates
the uncertainty of marginal calls.  A conjugate treatment is available at
negligible cost: with a Gamma prior $\lambda \sim \Gamma(a,b)$ (shape $a$, rate
$b$) fitted to a disjoint control region or to leave-one-locus-out global data,
$\lambda$ is integrated out rather than plugged in.  The integration must
include the observed-count factor $\lambda^n$ as well as the exposure, since
both carry $\lambda$:
\begin{equation}
    \log \int \lambda^{n}\exp\!\big(-2\lambda L_A^{\mathrm{eff}}\big)\,
      \pi(\lambda)\,d\lambda
    = \text{const} - (a+n)\log\!\big(b + 2L_A^{\mathrm{eff}}\big),
\end{equation}
so that the candidate-dependent comparison term is
\begin{equation}
    -(a+n)\,\log\frac{b+2L_1^{\mathrm{eff}}}{b+2L_2^{\mathrm{eff}}} .
\label{eq:gamma-term}
\end{equation}
This reduces to the plug-in $-2\bar\lambda(L_1^{\mathrm{eff}}-L_2^{\mathrm{eff}})$
with $\bar\lambda = a/b$ when the prior is tight relative to the exposure, and
widens as it loosens.  Omitting $\lambda^n$ --- integrating the exposure alone
--- is correct only when no reads were observed, and understates the term
substantially at realistic depth.  The shared flanks of the
recruited unit then serve as a genuine control contributing to $\pi(\lambda)$,
rather than as the informal cross-check described next.

\subsection{Cross-check on the flank}

The recruited unit contains sealed flanking sequence that both candidates share.
The observed read-start density there provides an independent estimate of
$\lambda$ that shares the local sequence context.  Agreement with the global
estimate is a validity check; disagreement is itself informative, indicating
locally depressed sampling.  Both values should be reported per locus, since a
marginal verdict cannot be interpreted without knowing which rate produced it.

\section{Reads explained by only one candidate}
\label{sec:truncation}

Because recruitment admits only reads matching recruited segments, reads
explicable by neither candidate never enter, and contamination requires no
model.  The delicate class is the opposite one: reads that place well in one
candidate and nowhere in the other.  These carry the most information, and they
are precisely the class recruitment \emph{selects for}.

For such a read, \eqref{eq:dllr} requires $\log\sum_p e_{jp}$ under the
candidate that cannot place it.  Evaluated literally this is the log-probability
of the best essentially random placement --- of order $10^4$ nats for a $15$\,kb
read, and scaling with read length.  A locus verdict would then be governed by
which reads happened to be long rather than by how much sequence is missing, and
the magnitude would rest on extrapolating a per-base error model tens of
thousands of bases into a range where it was never calibrated.

We therefore bound the contribution of each such read at a stated constant $C$.
Doing so changes the object being computed, and we distinguish the two
explicitly.  Write the uncapped ratio, over reads both candidates can place, as
\begin{equation}
    \Delta\mathrm{LLR}
    = \sum_{j\,\in\,\text{common}}\Bigg[
        \log\!\sum_{p\in\mathrm{slots}(A_1)}\! e^{(1)}_{jp}
      - \log\!\sum_{p\in\mathrm{slots}(A_2)}\! e^{(2)}_{jp}\Bigg]
      \;-\; 2\lambda\,\big(L_1^{\mathrm{eff}}-L_2^{\mathrm{eff}}\big),
\label{eq:llr-common}
\end{equation}
with $L_A^{\mathrm{eff}}$ from Proposition~\ref{prop:exposure}, and define the
reported statistic as
\begin{equation}
    \boxed{\;\Delta S \;=\; \Delta\mathrm{LLR} \;+\; C\,(m_1-m_2)\;}
\label{eq:final}
\end{equation}
where $m_1$ counts reads placeable only in $A_1$ and $m_2$ only in $A_2$.  Each
such read contributes through the count, not through the log-sum; the log-sums
run only over reads both candidates can place, which is the range in which the
error model was fitted.  This is Huberisation: bounding the influence of
observations lying where the model is least trustworthy.

\begin{remark}[$\Delta S$ is not a likelihood ratio]
\label{rem:not-lr}
Once the cap is imposed,
$\Delta S \neq \log\{\Prob(R\mid A_1)/\Prob(R\mid A_2)\}$.  It must not be
reported as a Bayes factor, converted to posterior odds, or described as
calibrated evidence.  It is a robust comparison score whose mapping to error
rates has to be established empirically, against simulated or held-out truth.
The placeability threshold and the value of $C$ are parameters of the
statistic: they are fixed in advance and swept as a sensitivity analysis, never
tuned per locus.  Under the ideal conditions of Section~\ref{sec:conditioning}
--- complete origin accounting, $\eta \to 0$, and no one-sided reads ---
$\Delta S$ reduces to $\Delta\mathrm{LLR}$, which is a likelihood ratio.
\end{remark}

\begin{proposition}[Extra sequence pays for itself, or does not]
\label{prop:count}
Let $A_1$ claim $\Delta$ bases that $A_2$ lacks, and let $q$ denote the
probability that a read originating in those bases is retained and uniquely
placeable.  If the sequence is real, $\mathbb{E}[m_1] = 2\lambda\Delta q$ and
\begin{equation}
    \mathbb{E}[\Delta S] = 2\lambda\Delta\,(Cq-1),
\end{equation}
positive whenever $C > 1/q$.  If the sequence is fabricated, no reads originate
from it, $m_1 = 0$, and $\Delta S = -2\lambda\Delta$.
\end{proposition}

The verdict direction is therefore stable over orders of magnitude in $C$
provided $C$ exceeds $1/q$; $C$ scales confidence only.  The detection
probability $q$ is bounded away from $1$ by the containment rule and the
placeability threshold, and must be estimated by simulation rather than
assumed.  Two readings of $m_1$ follow, and they should not be conflated.  If $q$ is
known and the Poisson model holds,
\begin{equation}
    \hat\Delta_1 = \frac{m_1}{2\lambda \hat q}
\end{equation}
is an unbiased estimator of the missing length $\Delta$.  If $q$ has not been
calibrated, then since $q \le 1$ and absent false positives,
$m_1/(2\lambda)$ has expectation $\Delta q \le \Delta$ and serves as a
conservative lower bound.  The first is the quantity to report once simulation
has supplied $\hat q$; the second is what may be claimed before that.  We recommend
reporting the triple $(m_1, m_2, L_1-L_2)$ alongside the score, since the counts
are robust to $C$, to $s$, and to the error model's tail behaviour entirely.

\section{The diploid case}
\label{sec:diploid}

Let a locus have haplotypes H1 and H2, each of length $\ell_{\mathrm{loc}}$,
differing at a fraction $d$ of sites.  Let $A_1$ carry both
($L_1 = 2\ell_{\mathrm{loc}}$) and $A_2$ carry only H1
($L_2 = \ell_{\mathrm{loc}}$).  Define the wrong-haplotype placement factor
\begin{equation}
    u = e^{-d\,\ell\,s} \in (0,1],
\end{equation}
the relative probability retained when a read is forced onto the other
haplotype.  Reads from H1 fit at both haplotypes in $A_1$ ($\propto 1+u$) and at
one in $A_2$ ($\propto 1$); reads from H2 fit at both in $A_1$ but must place on
H1 in $A_2$ ($\propto u$).  With $2\lambda\ell_{\mathrm{loc}} = n/2$,
\begin{equation}
    \Delta\mathrm{LLR}
    = \frac{n}{2}\Big[2\log(1+u) - \log u\Big] - \frac{n}{2}.
\end{equation}
The bracket has derivative $(u-1)/\big(u(1+u)\big) \le 0$ on $(0,1]$, so it is
minimised at $u=1$ with value $2\log 2 = 1.386 > 1$.  Hence
\begin{equation}
    \Delta\mathrm{LLR} \ge \left(\log 2 - \tfrac12\right) n \approx 0.193\,n > 0
\end{equation}
for all divergences, including zero --- but as an expectation, with $n$
replaced by its mean.  The realised score is random through the read counts of
each haplotype: at $d=0$ it reduces to
$N\log 2 - 2\lambda\ell_{\mathrm{loc}}$ with
$N \sim \text{Poisson}(4\lambda\ell_{\mathrm{loc}})$, whose failure
probabilities appear in Table~\ref{tab:power}.  At human-average
heterozygosity ($d=10^{-3}$, $\ell = 15$\,kb, $s=8$) and $n=400$ the expected
margin is of order $2\times10^4$ nats and the realised sign is never in
question; the finite-depth caveat binds only in the near-identical regime.  Haplotype loss is among the easiest events this method
detects, because the reads from the lost haplotype have nowhere to go.

A merged consensus haplotype rather than an absent one yields the same total
signal to leading order, $\approx n\,d\ell s/2$, since both read sets then
mismatch at half the heterozygous sites instead of one set mismatching at all of
them.  The two failure shapes need not be modelled separately.

\begin{remark}[Which term carries the haplotype signal]
\label{rem:diploid-class}
It is tempting to identify the H2 reads with the one-sided class $m_1$, but
that is wrong at realistic divergence.  Whether a read is placeable on the
surviving haplotype is decided by the pre-registered threshold, and a read of
length $\ell$ forced onto the wrong haplotype accrues about $d\ell$ extra
mismatches.  With an identity threshold $\theta$, it remains placeable while
$d \lesssim 1-\theta$: at $\theta = 0.99$ that is $d \lesssim 10^{-2}$, an
order of magnitude above human heterozygosity.  At $d = 10^{-3}$ and
$\ell = 15$\,kb the wrong-haplotype alignment retains an identity of about
$0.998$ and is comfortably retained.

The consequence is that at ordinary heterozygosity the H2 reads sit in the
\emph{common} class and the signal is carried by the emission ratio, each such
read contributing $\log\!\big((1+u)/u\big) \approx d\ell s \approx 120$ nats.
The one-sided classes carry haplotype signal only above the threshold's
divergence budget, and otherwise carry structural differences: with
$\theta = 0.99$ and $\ell = 15$\,kb, indels beyond roughly $135$\,bp.
Reporting $m_1$, $m_2$ and the emission term separately is therefore not
merely a diagnostic convenience --- the two record different kinds of event.
\end{remark}

\section{Assumptions and limits}

\begin{enumerate}
    \item \textbf{Uniform sampling within the unit.}  $\lambda$ is taken
          constant across the locus.  Where sequencing coverage is
          context-dependent --- notably in homopolymer-rich regions --- this is
          violated, and a correct assembly may be penalised for sequence that
          genuinely attracted few reads.  The flank cross-check of
          Section~\ref{sec:lambda} detects this; a context-modulated
          $\lambda(s)$ would mitigate it.
    \item \textbf{Haplotype-balanced coverage.}  Section~\ref{sec:diploid}
          assumes both haplotypes are sampled at rate $\lambda$.  Allele-specific
          depletion shades the verdict toward the collapsed candidate.
    \item \textbf{Read independence.}  Overlapping reads share genome, and HiFi
          errors are correlated within homopolymers and short tandem repeats.
          $\Delta\mathrm{LLR}$ therefore behaves as a Bayes factor scaled by
          an unknown factor and must not be read as calibrated odds, a
          restriction that applies a fortiori to $\Delta S$
          (Remark~\ref{rem:not-lr}).  Uncertainty
          should be obtained by a block bootstrap over genomic intervals rather
          than from the likelihood curvature.  The bootstrap addresses sampling
          variability only; the selection effect of
          Section~\ref{sec:validation} is separate and is not cured by it.
    \item \textbf{Origin enumeration.}  The sum over $p$ must not be truncated by
          an aligner's secondary-alignment limits, since that sum carries the
          copy-number signal.  Within a bounded unit, offsets should be
          enumerated exhaustively and scored by the forward algorithm.
    \item \textbf{Fixed error parameters.}  $\theta$ must be estimated once, from
          undisputed sequence, and held fixed across both candidates; fitting it
          per candidate lets a poorer assembly compensate with looser error
          rates.
\end{enumerate}

\section{Calibration and validation}
\label{sec:validation}

\subsection{The candidates were selected using these reads}

Both assemblies were built from the reads used to score them.  Treating them as
fixed hypotheses ignores the search that produced them: an assembler explores a
large space of alternatives and can encode sequencing error into consensus, so
a candidate's fit to its own input reads is optimistic in a way that no amount
of care in the likelihood repairs.  This is distinct from the dependence
between reads discussed below, and is not addressed by resampling: a block
bootstrap over a locus quantifies sampling variability, not selection.

Three responses, in decreasing order of strength.  A \emph{held-out library} ---
scoring on reads never supplied to either assembler --- removes the problem
outright.  \emph{Cross-fitting} approximates it: partition molecules into
folds, rebuild each candidate without one fold, and score it only on that fold.
Where neither is available, the honest description of $\Delta S$ is an internal
ranking statistic, useful for prioritising loci for curation, and not a
measure of evidence.

Two distinct forms of independence are easily conflated, and the validation
design turns on keeping them apart.

\emph{Independent truth} means an external standard against which a verdict can
be judged right or wrong.  A Q100 reference supplies this: it tells us whether
the method reached the correct conclusion, and it is what makes error rates
measurable at all.  Simulation supplies it by construction.

\emph{Independent reads} means scoring on data that did not participate in
building the candidates.  A Q100 reference does \emph{not} supply this: the
assemblies under comparison were still built from the same reads used to score
them, so $\Delta S$ remains optimistic even when the benchmark says the verdict
was right.  What an external truth set buys is an unbiased estimate of the
\emph{decision} accuracy; what it does not buy is an unbiased \emph{score}.

The practical consequence: HG002-against-Q100 is a sound basis for calibrating
the mapping from $\Delta S$ to error rates, provided the calibration is
reported as accuracy of the decision rather than as evidence carried by the
score, and provided the same reads were used to build both candidates so the
optimism is at least symmetric.  A survey ranking production assemblies against
one another with neither external truth nor held-out reads has neither
property, and its output is a prioritisation for curation and nothing more.

\subsection{What must be calibrated}

Three quantities in this document are not determined by the model and must be
established empirically, all by simulation against known truth:

\begin{enumerate}
    \item the mapping from $\Delta S$ to error rates, which
          Remark~\ref{rem:not-lr} defers and which is the only thing that makes
          the score interpretable;
    \item the detection probability $q$ of Proposition~\ref{prop:count},
          without which $m_1$ bounds missing length but does not estimate it;
    \item the gap and insertion parameters of Section~\ref{sec:context}, fitted
          on undisputed sequence and held fixed.
\end{enumerate}

Sensitivity should be reported over the placeability threshold, the cap $C$,
and $\lambda$, since each is a choice rather than a measurement.  Because the
propositions of Sections~\ref{sec:conditioning} and~\ref{sec:diploid} concern
expectations, calibration must report realised error rates at finite depth ---
Table~\ref{tab:power} gives the floor those rates cannot beat.

\section{Summary}

The verdict for a locus is $\Delta S$ of \eqref{eq:final}, comprising three
terms: a
likelihood difference over reads both candidates explain, computed by a
free-start forward pass with quality-conditioned and context-dependent
parameters; a
bounded count difference over reads only one candidate explains; and a linear
exposure penalty $2\lambda(L_1^{\mathrm{eff}}-L_2^{\mathrm{eff}})$ with
$\lambda$ estimated from the k-mer spectrum of the reads.  The exposure penalty
and the count are two halves of one question --- whether claimed sequence pays
for itself in reads --- expressed in the same units.

Three qualifications travel with every verdict.  Its likelihood-ratio component
is local only to the tolerance $\eta$ of Proposition~\ref{prop:locality}, which
is computed rather than assumed and reported alongside the score; the capped
term is outside that guarantee, which is why $m_1$ and $m_2$ are reported
separately (Remark~\ref{rem:eta-scope}).  It is a robust statistic rather
than calibrated evidence, reducing to a likelihood ratio only under the ideal
conditions of Remark~\ref{rem:not-lr}.  And the propositions establish the sign
of an expectation, not of every realisation: at low depth the realised verdict
can invert, with the floor given in Table~\ref{tab:power}.

Stated at the width the derivations support: for fixed candidates and an
explicitly modelled read-selection process, the method supplies an
exposure-calibrated comparison score, which under ideal closure, complete origin
accounting and a calibrated mark likelihood reduces to a likelihood ratio whose
copy-number signal survives because local coverage is never profiled away.

\end{document}
